\documentclass[a4paper]{article}
\usepackage{amsfonts}
\usepackage{amsmath}
\usepackage{amssymb}
\usepackage{graphicx}

\begin{document}

\noindent \large Auteur: Rob Willemsen \\
\noindent \large Studierichting: Informatica\\
\noindent \large Oefeningenreeks 1C nr. 12.3\\

\medskip
\normalsize

Gegeven is de breuk $z = \dfrac{\sqrt{2}}{1+i}$. We zoeken de goniometrische vorm door eerst teller en noemer afzonderlijk om te zetten. \\

\textbf{Teller $z_1 = \sqrt{2}$:} \\

$r_1 = \sqrt{(\sqrt{2})^2 + 0^2} = \sqrt{2}$ \\

$\alpha_1 = 0 \Rightarrow z_1 = \sqrt{2} \operatorname{cis} 0$ \\

\textbf{Noemer $z_2 = 1+i$:} \\

$r_2 = \sqrt{1^2 + 1^2} = \sqrt{2}$ \\

$\tan \alpha_2 = \dfrac{1}{1} = 1 \Rightarrow \alpha_2 = \dfrac{\pi}{4} \Rightarrow z_2 = \sqrt{2} \operatorname{cis} \dfrac{\pi}{4}$ \\

\textbf{Quotient bepalen:} \\

$z = \dfrac{z_1}{z_2} = \dfrac{\sqrt{2} \operatorname{cis} 0}{\sqrt{2} \operatorname{cis} \dfrac{\pi}{4}}$ \\

$z = \dfrac{\sqrt{2}}{\sqrt{2}} \operatorname{cis} \left(0 - \dfrac{\pi}{4}\right)$ \\

$z = 1 \operatorname{cis} \left(-\dfrac{\pi}{4}\right)$ \\

\textbf{Antwoord:} $z = \operatorname{cis} \left(-\dfrac{\pi}{4}\right)$ \\

\begin{thebibliography}{999}
%\bibitem{a} Referentie
\end{thebibliography}

\end{document}